% Chapter 8 -- Learning More About FLAIR (Part I, Learning; the closing
% chapter of the course).  Sources (Law 1): os/flair/finder_menu.c (the
% Finder's Help menu: About Help and Show Balloons always dimmed);
% os/milton/kmain.c (HELLO and NOTES launched at start-up; finder_surface_
% open: an .EXE icon is launched); the data volume as minted by the
% Makefile (FLAIR_FIXTURES: README.TXT at the top level, the applications in
% APPS); os/apps/itext/itext.c (Open...: README.TXT read into the window).
% Model: Macintosh User's Guide 1991 Ch. 7 "Learning More About Your
% Macintosh" (apple-macintosh.md Sec 8.2).
% GATES: test-flair-desktop test-finder-menu test-flair-itext test-stdfile test-flair-app-launch test-pubs-flair-screens
\Gates{test-flair-desktop, test-finder-menu, test-flair-itext,
  test-stdfile, test-flair-app-launch, test-pubs-flair-screens}
\chapter{Learning More About FLAIR}
\label{ch:learning}

You have now done everything you need for everyday work with FLAIR: you
can start and stop the computer, use the mouse and the menus, arrange
windows, organize your files, use an application, write and save a
document, and throw away what you no longer need.  This chapter tells you
where to go from here.

\section{Reading the README file}
\index{README.TXT}
The data disk holds a short text file, \file{README.TXT}, with notes about
the disk.  Read it with InitechText.
\begin{TaskSteps}
\Step Open InitechText (Chapter~\ref{ch:applications}).
\Step Choose \menuitem{Open...} from the \menuname{File} menu.
\Step Double-click \file{README.TXT} in the list.
\Result{The file opens in InitechText's window.}
\end{TaskSteps}

\FlairScreen[label=fig:readme]{l_readme}{README.TXT, opened in InitechText}

The notes refer to the InitechDOS command prompt as the \textit{DOS
prompt}, and to folders as \textit{directories}: Chapter~\ref{ch:dos}
explains how FLAIR and InitechDOS name the same things.  When you have read
the file, quit InitechText without saving.

\section{Practising}
\index{practising}
The \texttt{HELLO} and \texttt{NOTES} windows that FLAIR opens when it
starts are there for you to practise on.  Move them, change their size,
zoom and collapse them, pull down their menus, and bring each to the front
in turn.  Nothing you do to them changes anything on the disk.  When you
close one with its close box, it stays closed until you next start FLAIR.

To practise organizing files, make a new folder in the \texttt{Drive A
Files} window, move \file{README.TXT} into it and back out again, and then
throw the empty folder away and empty the Trash.

\section{Using this book as a reference}
\index{reference}
Part~II of this book is for looking things up, in any order:
\begin{itemize}
\item Chapter~\ref{ch:menus} describes every command in every menu.
\item Chapter~\ref{ch:windowsref} describes every kind of window, and
  Chapter~\ref{ch:alerts} every alert and dialog box.
\item Chapter~\ref{ch:disks} describes the disks and what is on them, and
  Chapter~\ref{ch:adapting} what you can arrange to suit your work.
\item Chapter~\ref{ch:trouble} tells you what to do when something does
  not work as you expect.
\end{itemize}
The Summary at the end of Part~II gathers the common tasks in short tables,
and the index lists every subject the book covers.

\section{The Help menu}
\index{Help menu}\index{on-line help}
The Finder's menu bar ends with a \menuname{Help} menu.  Its two commands,
\menuitem{About Help} and \menuitem{Show Balloons}, are dimmed: this
version of FLAIR has no help on the screen.  Use this book, the Summary and
the \doctitle{FLAIR Quick Reference} card instead.

\section{The other applications}
\index{applications!other}
The \file{APPS} folder holds more applications than the ones this book
describes.  Each opens from the desktop in the same way: double-click its
icon.
\begin{center}\small
\begin{tabular}{@{}p{24mm}p{75mm}@{}}
\textbf{Icon} & \textbf{Where it is described}\\\hline
\file{TEXT.EXE} & InitechText: Chapter~\ref{ch:documents}, and the
\doctitle{FLAIR Accessories Guide}\\
\file{CALCULAT.EXE} & the Calculator: the \doctitle{FLAIR Accessories
Guide}\\
\file{CALC.EXE} & InitechCalc: the \doctitle{FLAIR Accessories Guide}\\
\file{PAINT.EXE} & InitechPaint: the \doctitle{FLAIR Accessories Guide}\\
\file{CTENANT.EXE}, \file{TENANTFX.EXE} & the sample applications: the
\doctitle{FLAIR Accessories Guide}\\
\file{123.EXE} & Initech 123: its own books\\
\file{WORD.EXE} & InitechWord: its own books\\
\end{tabular}
\end{center}
Remember that one application from the disk can be open at a time: quit one
before you open the next (Chapter~\ref{ch:multi}).

\section{The other books}
\begin{itemize}
\item The \doctitle{FLAIR Accessories Guide} describes each of the small
  applications that come with FLAIR.
\item The \doctitle{FLAIR Quick Reference} card shows every menu and every
  keyboard equivalent on one sheet.  Keep it beside the computer.
\item The \doctitle{InitechDOS User's Guide} and \doctitle{InitechDOS
  Reference} describe the command prompt you reach with \menuitem{Exit to
  InitechDOS}.
\end{itemize}
Appendix~D lists every book with its part number, and tells you how to
reach the Initech Help Desk.
